\documentclass[11pt]{article}
\usepackage[margin=1in]{geometry}
\usepackage{amsmath,amssymb,amsthm,mathtools}
\usepackage{graphicx}
\usepackage{microtype}
\usepackage{xcolor}
\usepackage[colorlinks=true,linkcolor=blue!55!black,citecolor=blue!55!black,urlcolor=blue!55!black]{hyperref}
\usepackage{enumitem}

\newtheorem{theorem}{Theorem}
\newtheorem{lemma}[theorem]{Lemma}
\newcommand{\cS}{\mathcal S}
\newcommand{\cM}{\mathcal M}

\title{The Type-I Hypergraph Hanner Conjecture}
\author{Candidate full solution for arXiv:0903.3237\\Model: GPT5.6}
\date{19 July 2026}

\begin{document}
\maketitle

\begin{center}
\fbox{\parbox{0.91\textwidth}{\textbf{Status: full solution, likely valid.}
Every semi-norming Type-I hypergraph pair with parameter $s\ge2$ satisfies
the conjectured $|H|$-Hanner inequality.  This proves Conjecture~3.15 of
Hatami~\cite{Hatami}; non-factorizability is not needed.}}
\end{center}

\section{Source problem}

Let $H=(\alpha,\beta)$ be a finite hypergraph pair.  In Hatami's notation,
$H$ is of Type~I with parameter $s$ if, on its support $\cS$,
\[
 \alpha(\omega)=\beta(\omega)=\frac{s}{2}.
\]
If $m=|\cS|$, then $|H|=ms$.  For a function $f$ on the underlying product
measure space, the associated hypergraph norm satisfies
\begin{equation}\label{eq:type-i}
 \|f\|_H^{ms}
 =\int \prod_{\omega\in\cS}|f(\omega(x))|^s\,dx.
\end{equation}
Conjecture~3.15 asks whether, for $s\ge2$,
\begin{equation}\label{eq:goal}
 \|f+g\|_H^{|H|}+\|f-g\|_H^{|H|}
 \le (\|f\|_H+\|g\|_H)^{|H|}
      +|\|f\|_H-\|g\|_H|^{|H|}.
\end{equation}

\begin{figure}[ht]
  \centering
  \includegraphics[width=0.98\textwidth]{figures/open_problem_crop.png}
  \caption{Conjecture~3.15 on page 25 of the source paper.}
\end{figure}

The source proves the corresponding inequality on $\ell_H$ when $s$ is an
even integer by expanding the $s$th powers directly.  The point below is to
perform only the finite expansion indexed by the $m$ support elements, leaving
the real exponent $s$ inside an ordinary $L_s$ norm; the argument works on an
arbitrary underlying measure space.

\section{Two elementary ingredients}

We use Hatami's generalized H\"older inequality \cite[Lemma~2.10]{Hatami}.
In the special form needed here, if $A\subseteq\cS$ and
$a=\|f\|_H$, $b=\|g\|_H$, then
\begin{equation}\label{eq:mixed-holder}
 \int \prod_{\omega\notin A}|f(\omega(x))|^s
       \prod_{\omega\in A}|g(\omega(x))|^s\,dx
 \le a^{s(m-|A|)}b^{s|A|}.
\end{equation}
Indeed, decompose $H$ into the two nonnegative hypergraph pairs supported on
$\cS\setminus A$ and $A$, and apply the nonnegative-function clause of that
lemma to $|f|$ and $|g|$.

We will also need the following small monotonicity observation.

\begin{lemma}\label{lem:monotone}
For $s\ge1$, the function
\[
 \Phi_s(u,v)=(u+v)^s+|u-v|^s
\]
is nondecreasing in each variable on $[0,\infty)^2$.
\end{lemma}

\begin{proof}
On the region $u\ge v$, its partial derivatives are
\[
 s\big((u+v)^{s-1}+(u-v)^{s-1}\big),\qquad
 s\big((u+v)^{s-1}-(u-v)^{s-1}\big),
\]
and both are nonnegative.  The other region follows by symmetry.
\end{proof}

\section{Idea of the proof}

The obstruction in the source proof is the temptation to expand the real
$s$th power, which is possible only when $s$ is an even integer.  Instead we
factor the Type-I integrand as the ordinary $L_s$ norm of the product over the
$m$ support elements.  That product has an integer-degree expansion for every
real $s$: its even and odd terms form two functions $E$ and $O$.  The ordinary
$L_s$ Hanner inequality controls $E+O$ and $E-O$, while Hatami's generalized
H\"older lemma controls each mixed monomial in $E$ and $O$.  The two resulting
bounds are precisely the even and odd binomial sums; coordinatewise
monotonicity of the two-variable Hanner expression then recombines them into
the conjectured right-hand side.

\section{Full solution}

\begin{theorem}\label{thm:main}
Let $H$ be a semi-norming hypergraph pair of Type~I with parameter $s\ge2$.
Then every $L_H$ space satisfies the $|H|$-Hanner inequality
\eqref{eq:goal}.  In particular, Conjecture~3.15 of \cite{Hatami} holds.
\end{theorem}

\begin{proof}
Fix $f,g\in L_H$, and abbreviate $a=\|f\|_H$, $b=\|g\|_H$.
For $A\subseteq\cS$, put
\[
 M_A(x)=\prod_{\omega\notin A}f(\omega(x))
        \prod_{\omega\in A}g(\omega(x)).
\]
Separate the expansion into its even and odd parts:
\[
 E=\sum_{|A|\ \mathrm{even}}M_A,
 \qquad
 O=\sum_{|A|\ \mathrm{odd}}M_A.
\]
Ordinary distributivity gives
\begin{equation}\label{eq:evenodd}
 E+O=\prod_{\omega\in\cS}(f(\omega(x))+g(\omega(x))),
 \qquad
 E-O=\prod_{\omega\in\cS}(f(\omega(x))-g(\omega(x))).
\end{equation}

Apply the classical Hanner inequality in the ordinary complex $L_s$ space
\cite{Hanner}:
\begin{equation}\label{eq:ordinary-hanner}
 \|E+O\|_s^s+\|E-O\|_s^s
 \le \Phi_s(\|E\|_s,\|O\|_s).
\end{equation}
By \eqref{eq:mixed-holder},
\[
 \|M_A\|_s\le a^{m-|A|}b^{|A|}.
\]
The triangle inequality in $L_s$ therefore yields
\begin{align*}
 \|E\|_s&\le E_0:=\sum_{k\ \mathrm{even}}\binom{m}{k}a^{m-k}b^k,\\
 \|O\|_s&\le O_0:=\sum_{k\ \mathrm{odd}}\binom{m}{k}a^{m-k}b^k.
\end{align*}
Lemma~\ref{lem:monotone} permits these two upper bounds to be substituted
simultaneously into \eqref{eq:ordinary-hanner}.  The even and odd binomial
identities give
\[
 E_0+O_0=(a+b)^m,
 \qquad E_0-O_0=(a-b)^m.
\]
Consequently,
\begin{equation}\label{eq:finish}
 \|E+O\|_s^s+\|E-O\|_s^s
 \le (a+b)^{ms}+|a-b|^{ms}.
\end{equation}
Finally, \eqref{eq:type-i} and \eqref{eq:evenodd} identify the left-hand
side of \eqref{eq:finish} with
$\|f+g\|_H^{ms}+\|f-g\|_H^{ms}$.  Since $ms=|H|$, this is exactly
\eqref{eq:goal}.
\end{proof}

\section{Checks and scope}

\begin{enumerate}[leftmargin=2em]
\item The proof never expands a noninteger power.  It expands a product with
integer degree $m$ and then invokes the classical $L_s$ inequality.
\item Complex-valued functions cause no change: \eqref{eq:evenodd} is
algebraic, the mixed H\"older estimate is applied to absolute values, and the
classical Hanner inequality holds in complex $L_s$.
\item Non-factorizability is absent from the argument.  Thus the theorem is
slightly stronger than the source conjecture.
\item The restriction $s\ge2$ occurs only in the direction of the ordinary
$L_s$ Hanner inequality.  For $s<2$ its direction reverses, consistent with
the obstruction recorded immediately before Conjecture~3.15.
\end{enumerate}

For a numerical diagnostic, the accompanying script tests the first
non-even example: the Type-I norm generated by the four-cycle with $s=3$.
Random search and differential evolution found no positive Hanner defect;
the optimized value was within $3.1\cdot10^{-15}$ of equality.  This check is
not used in the proof.

\section{Novelty note}

The run's registry, solution, attempt, and proof-gap indexes contained no
entry for arXiv:0903.3237.  Exact searches for the source conjecture,
``Type-I hypergraph pair,'' and ``Hanner inequality'' found Hatami's paper
and thesis, but no later proof.  Nearby literature on complex $L_p$ Hanner
inequalities and Schatten norms does not state the even/odd hypergraph
argument above.  This was a bounded search, not a complete citation audit.

\begin{thebibliography}{9}
\bibitem{Hatami}
H. Hatami,
\emph{On generalizations of Gowers norms and their geometry},
arXiv:0903.3237 (2009), especially Lemma~2.10 and Conjecture~3.15.

\bibitem{Hanner}
O. Hanner,
\emph{On the uniform convexity of $L^p$ and $\ell^p$},
Arkiv f\"or Matematik 3 (1956), 239--244.
\end{thebibliography}

\end{document}
